%%%%%%
% Metadata
%%%%%%

\chapter{Introduction}
\label{chapter:Intro}
\section{Background and Motivation}

Statistical hypothesis testing is one of the most commonly used methods in scientific research. In empirical research, researchers aim to assess whether an observed effect represents a real relationship or whether it may have occurred due to random variation.

To answer this question, researchers define two competing statements. The first is the null hypothesis, denoted by $H_0$, which usually represents "no effect" or "no difference". The second is the alternative hypothesis, denoted by $H_1$, which represents the presence of an effect.

After collecting data, researchers calculate a value known as the \emph{$p$-value}. \cite{Fisher1925} introduced the $p$-value as a way to measure how compatible the observed data are with the null hypothesis. In simple terms, the $p$-value answers the following question:

\begin{quote}
	If the null hypothesis were true, how likely would it be to observe data like the ones we obtained?
\end{quote}

The \textit{$p$-value} is the probability of observing a test statistic at least as extreme as the observed value, assuming that the null hypothesis is true. It measures how compatible the observed data are with the null hypothesis.

More formally, if $T$ is a test statistic and $t_{\text{obs}}$ is its observed value, the $p$-value can be written as

\begin{equation}
	p = P(T \geq t_{\text{obs}} \mid H_0)
	\end{equation}



To make a decision, the $p$-value is compared to a pre chosen threshold called the significance level, denoted by $\alpha$. According to the framework developed by \cite{NeymanPearson1933}, the null hypothesis is rejected if
\begin{equation}
	p \leq \alpha
\end{equation}

The significance level $\alpha$ represents the probability of making a Type I error, which means rejecting the null hypothesis even though it is true \cite{NeymanPearson1933}. Formally,

\begin{equation}
	\alpha = P(\text{Reject } H_0 \mid H_0 \text{ is true})
\end{equation}
This procedure, known as the Neyman–Pearson approach to statistical inference, allows researchers to make binary claims while controlling the error rate\cite{Lakens2021}
In practice, researchers often choose $\alpha = 0.05$. Other common values are $0.10$ and $0.01$. The value $0.05$ has become the standard in many scientific fields. However, \cite{Fisher1925} originally suggested $0.05$ as a convenient reference point, not as a universal rule. Later research has shown that the dominance of $0.05$ is mainly historical and conventional rather than theoretically optimal \cite{CowlesDavis1982}.

At first glance, this decision rule seems clear and simple. If the $p$-value is smaller than $0.05$, the result is called "statistically significant". If it is larger, it is not.

However, problems arise when the $p$-value is very close to the threshold. For example, if $p = 0.049$, the null hypothesis is rejected. If $p = 0.051$, it is not rejected. The numerical difference is very small, but the conclusion changes completely. In practice, researchers sometimes describe such results as "almost significant" or "marginally significant", even though hypothesis testing formally allows only two decisions\cite{Pritschet2016}.

This discomfort has become more serious in recent years. Large replication projects have shown that many statistically significant findings cannot be reproduced in independent studies. For example, the Open Science Collaboration \cite{OpenScience2015} reported that a substantial proportion of findings in psychology failed to replicate . In addition, the ASA (American Statistical Association) warned that $p$-values should not be used mechanically or interpreted without context \cite{WassersteinLazar2016}.

These developments suggest that the routine use of a fixed significance level, especially $\alpha = 0.05$, may not always be appropriate.  More recent work confirms that replication failure remains a persistent challenge across disciplines  and that the validity of initial findings is systematically overestimated when significance thresholds are applied uniformly \cite{Nosek22}.

\section{Research Problem}

Even though many researchers are aware of the limitations of fixed significance levels, the value $\alpha = 0.05$ is still used almost automatically in most empirical studies. This creates several concerns.

First, studies differ in important ways. Some studies have small sample sizes, while others include thousands of observations. Some studies aim to detect very small effects, while others focus on large effects. Despite these differences, the same significance level is often applied.

Additionally, fixing $\alpha$ without accounting for sample size 
or expected effect size can create an imbalance between false 
positive and false negative error rates, a relationship explored 
in detail in Chapter~\ref{chapter:history}.

Statistical power refers to the probability of rejecting a false null hypothesis\cite{Aberson2015}. Power depends on three main factors: the significance level $\alpha$, the sample size, and the true effect size. \cite{Cohen1988} emphasized that these quantities are closely related and should be considered together when planning a study. Fixing $\alpha$ without considering sample size or expected effect size may therefore lead to an imbalance between false positive and false negative decisions.

Third, the strict division between "significant" and "not significant" can create misleading interpretations when $p$-values are close to the threshold. Two results that are almost identical numerically can lead to completely different conclusions simply because one falls slightly below $0.05$ and the other slightly above.

Although the statistical literature contains proposals to modify or lower the conventional threshold, such as recommending $\alpha = 0.005$ for certain claims \cite{Benjamin2018}, these proposals still rely on fixed values. To the best of the author's knowledge, no widely adopted method exists that systematically adapts the significance level to the specific characteristics of an individual study.

This gap forms the starting point of the present thesis.
\section{Research Objective}

The main objective of this thesis is to develop and evaluate a systematic approach for choosing the significance level $\alpha$. Instead of treating $\alpha$ as a fixed constant, this thesis proposes to view it as a flexible value that can depend on study characteristics. In particular, the proposed approach uses information from past replication studies to inform the choice of $\alpha$. Study characteristics such as sample size and effect size are used to generate a data based baseline recommendation using a machine learning model.

In a second step, this baseline value can be adjusted based on contextual factors, such as the researcher’s tolerance for risk or the number of hypotheses being tested.

The goal is not to replace statistical reasoning or to claim that there is one universally correct significance level. Instead, the goal is to provide a structured and transparent recommendation that helps researchers make more informed decisions, especially in situations where the $p$-value is close to the threshold.

The central research question of this thesis is:

\begin{quote}
	\textit{Can a data driven and context sensitive approach 
		improve the choice of the significance level compared to 
		the conventional fixed value of $\alpha = 0.05$?}
\end{quote}


\section{Structure of the Thesis}

The remainder of this thesis is organized as follows:

\begin{itemize}
	
	\item Chapter \ref{chapter:history} introduces the historical foundations of statistical hypothesis testing. It explains how significance testing developed and discusses the main ideas of Fisher and the Neyman–Pearson framework. It also describes how the commonly used significance level $\alpha = 0.05$ became widely adopted in scientific research.
	
	\item Chapter \ref{chapter:existing approaches} reviews previous research on how to choose a significance level. It discusses several approaches from the literature, including fixed thresholds, power based methods, and decision theoretic approaches. The limitations of these methods are also highlighted.
	
	\item Chapter \ref{chapter:development} explains the development of the proposed Alpha Recommendation System. It describes how the dataset is constructed, how optimal alpha values are defined using replication studies, and how a machine learning model is used to predict $\alpha_{base}$. The chapter also explains the context based adjustment step that modifies the $\alpha_{base}$ based on study characteristics.
	
	\item Chapter \ref{chapter:deplyment} presents the deployment of the system. It explains how the trained machine learning model is integrated into an interactive application using Streamlit and how users can interact with the system to obtain an alpha recommendation.
	
	\item Chapter \ref{chapter:testing} evaluates the performance of the proposed system. The recommended alpha values are tested using selected studies and compared with decisions based on the conventional significance level.
	
	\item Chapter \ref{chapter:Conclusion} concludes the thesis by summarizing the main findings, discussing the limitations of the approach, and suggesting possible directions for future research.
	
\end{itemize}


